\maketitle
\tableofcontents

% ============================================================
\chapter*{Doc.~314 -- The Lattice in Hilbert Space}
\addcontentsline{toc}{chapter}{Doc.~314 -- The Lattice in Hilbert Space}
% ============================================================

\section*{What this is about}

FFGFT computes geometrically on T$^4$ with the D4 lattice and
quantum-mechanically on $\mathcal{H} = L^2(\text{T}^4)\otimes
\mathbb{C}^3$ (Docs.~230/282). The transition raises an obvious
question: \emph{do the lattices become flat?} Does anything remain of
kissing number~24, packing density and triality, or is every lattice
the same in Hilbert space?

The answer: the torus is flat anyway -- curvature was never the
difference -- and the lattice structure is not lost but
\textbf{migrates into the spectrum and into the fibre}. This document
computes that through (three scripts, Ch.~J) and finds:

\begin{enumerate}
  \item \textbf{The spectrum separates what curvature does not
        [K]:} Z$^4$ starts at $|k|^2 = 1$ with degeneracy~8, D4 at
        $|k|^2 = 2$ with degeneracy~24; odd norms are entirely
        absent for D4. The first excitation carries exactly
        $24 = 12+12$ -- FCC half and winding half separately
        addressable.
  \item \textbf{Triality is built in, not attached [K]:}
        $|\text{Aut}(D4)| = 1152 = |W(F_4)|$, the quotient to
        $W(D4)$ is $6 = |S_3|$; among the 80 order-3 elements, 16
        generate exactly the T$^4/Z_3$ orbifold with
        $|\det(1-A)| = 9$ fixed points and act on every mode orbit
        as a $\mathbb{C}^3$ circulant.
  \item \textbf{A negative finding with theorem character [K]:} $5$
        does not divide $1152 = 2^7\cdot3^2$, so $\text{Aut}(D4)$
        contains \emph{no} element of order~5. The fivefold rotation
        from which $\theta = 2/9$ arises (Doc.~293) is in principle
        incompatible with the lattice -- the phase belongs to a
        different layer.
  \item \textbf{The doubling is nevertheless prepared [K]:}
        $24 = 8\times3 = 4\times6$; four disjoint six-blocks with
        exactly the $C_3$ content that $3\oplus3'$ requires.
  \item \textbf{The reciprocity is mode-wise exact -- as a
        \emph{place} relation [K]:} $\lambda_4\cdot m = 2\pi$
        identically for every mode (de Broglie along the mass
        direction; winding number $=$ mass number). The
        \emph{temporal} reading carries $1/\gamma$, a mass mixture
        a Jensen excess $\ge1$: the relation becomes exact only
        for the resting, sharp mode.
  \item \textbf{The perturbation calculation is carried out [K]:}
        multiplet sizes follow the irrep dimensions of the
        respected symmetry, not the orbits -- $9+8+4+2+1$ in the
        fully symmetric case, $8\times1+8\times2$ for $Z_3$,
        generically 24 individual levels.
  \item \textbf{The closure fork is spectrally invisible [K]:} in
        the rolled-up state there are no fractal corrections --
        winding numbers are topological, the defect $d = 100\xi$
        accumulates only when unrolled. Cases A/B/C of
        Docs.~313/295 leave the mode spectrum untouched, so
        Chapters~E--H are case-independent. The same boundary
        settles the fractal question (Ch.~J): $K_\text{frak}$ is
        cumulative (Doc.~133) and acts only through the scale
        bridge; locally only $D_f^\text{space} = 3-\xi$ acts, with
        $6.7\times10^{-5}$.
\end{enumerate}

Not part of the A-series; status: working document. References:
Docs.~230/282 ($\mathcal{H} = L^2(\text{T}^4)\otimes\mathbb{C}^3$),
133 (derivation of $K_\text{frak}$: cumulative, RG run), 285
(doubling $6 = 3\oplus3'$), 291 (channel analysis), 293
($\theta = 2/9$ icosahedral), 295/313 (closure fork), 306 (native time--energy
reciprocity), 311 (Four onto Three, $24 = 12+12$).

% ============================================================
\section*{Reader's guide: the thread of the argument}
% ============================================================

\textbf{The question.} In Hilbert space the lattice becomes the
index set of the Fourier modes. It is then no longer a sphere
packing -- in that sense every lattice becomes \enquote{flat}. The
question is whether the information disappears with it.

\textbf{The answer in one line.} \emph{Geometry $\to$ spectrum $+$
fibre.} The degeneracies of the Laplacian are the theta series of
the lattice; the kissing number becomes the degeneracy of the first
excitation; triality becomes the orbifold and fibre structure.

\textbf{What follows.} First, two flat tori are distinguishable by
their spectra (Ch.~B). Second, the $\tilde T\cdot m = 1$ structure
becomes spectrally readable when $R_4 \ne R_{123}$, and the level
ordering changes at exactly two root radii (Ch.~C). Third, the
$Z_3$ need not be imposed -- it is a lattice symmetry (Ch.~E).

\textbf{The boundary, drawn cleanly.} Fivefold symmetry is
impossible on the lattice (Ch.~F). This does not hit the doubling,
because that lives on the \emph{fibre} and the interface of the two
levels is exactly $C_3$ -- occupied by triality (Ch.~G). And the
container structure holds on \emph{every} shell, so it cannot
distinguish one (Ch.~H) -- the same pattern-versus-occupation logic
as in Doc.~293.

\textbf{The caveat.} The computation uses the smooth Laplacian,
i.e.\ $D_f = 4$ exactly. All numbers are therefore \emph{bare};
Ch.~I separates what that costs: counts and ratios stay invariant,
$D_f^\text{space} = 3-\xi$ is \emph{local} and acts with
$6.7\times10^{-5}$, while $K_\text{frak} = 1-100\xi$ is
\emph{cumulative} (RG run across 17 orders of magnitude) and acts
only through the scale bridge. Bare are therefore only the SI
absolute values ($-1.33\,\%$ via the anchor); counts, ratios and
containers are not. The degeneracy splitting under perturbation is
computed -- it follows the irrep dimensions of the perturbation
symmetry, not the orbits, and lies at $\sim10^{-5}$. Ch.~D2 adds
the boundary between rolled-up and unrolled states, which is why
the closure fork (Docs.~313/295) stays spectrally invisible.

\medskip
\noindent\textbf{In one sentence:} in Hilbert space the lattice
becomes flat but not inconsequential -- what geometry carried is
carried by spectrum and fibre, and the division of labour between
them is exactly determinable, as long as the smooth approximation is
booked as such.

% ============================================================
\section*{Notation}
% ============================================================

The table explains the symbols of this document. Quantities
appearing only once are introduced there.

\textbf{Note on notation.} In lattice theory $\Lambda$ is the usual
symbol for a lattice (Conway/Sloane). In the FFGFT corpus, however,
$\Lambda$ is taken -- as the cosmological constant, in
$\Lambda$CDM, and in particular as the closure scale
$\Lambda^{*} = 1/R_H^2$ (Doc.~312). This document therefore uses
$\Gamma$ throughout for the lattice and $\Gamma^{*}$ for the dual
lattice. Likewise $L^{*}$ is \emph{not} used here, since in
Doc.~312 it denotes the cycle length; the torus edge lengths are
called $R_{123}$ and $R_4$.

\begin{center}
\renewcommand{\arraystretch}{1.3}
{\small\setlength\tabcolsep{5pt}
\begin{longtable}{lp{9.6cm}}
\toprule
\textbf{Symbol} & \textbf{Meaning} \\
\midrule
\endfirsthead
\toprule
\textbf{Symbol} & \textbf{Meaning} \\
\midrule
\endhead
\bottomrule
\endfoot
\multicolumn{2}{l}{\textit{Lattice and torus}} \\
$\Gamma$ & lattice in $\mathbb{R}^4$; defines the torus
  T$^4 = \mathbb{R}^4/\Gamma$ (instead of the $\Lambda$ customary
  in the lattice literature, see note on notation) \\
$\Gamma^{*}$ & dual lattice; carries the admissible wave vectors
  of the Fourier modes \\
$D4$ & root lattice $D_4$: all integer $k$ with $\sum k_i$ even;
  kissing number 24 \\
$\text{Z}^4$ & cubic lattice $\mathbb{Z}^4$; kissing number 8,
  exactly self-dual \\
$D4^{*}$ & dual D4 lattice; again a D4 lattice up to scaling \\
$R_{123}$, $R_4$ & radii of the three spatial directions and of the
  fourth direction \\
$r$ & radius ratio $r = R_4/R_{123}$ (dimensionless) \\
\midrule
\multicolumn{2}{l}{\textit{Spectrum}} \\
$k$ & wave vector (mode index), $k \in \Gamma^{*}$ \\
$k_4$ & fourth component of $k$; $k_4 = 0$ separates the FCC half
  from the winding half \\
$|k|^2$ & norm; also the eigenvalue of the Laplacian on the
  isotropic torus \\
$E(k)$ & eigenvalue under anisotropy,
  $E = k_1^2+k_2^2+k_3^2+(k_4/r)^2$ \\
$N(n)$ & degeneracy of the shell $|k|^2 = n$; coefficient of the
  theta series \\
$\theta_{D4}$ & theta series of the D4 lattice,
  $\theta_{D4} = (\theta_3^4+\theta_4^4)/2$ \\
$\theta(t)$ & heat-kernel theta function,
  $\theta(t) = \sum_k e^{-t|k|^2}$ \\
$\zeta_M(s)$ & Epstein zeta function of the lattice $M$;
  $s = -1/2$ is the Casimir point \\
\midrule
\multicolumn{2}{l}{\textit{Symmetries}} \\
$\text{Aut}(\Gamma)$ & automorphism group of the lattice
  (length-preserving self-maps) \\
$W(F_4)$, $W(D4)$ & Weyl groups; $|W(F_4)| = 1152$,
  $|W(D4)| = 192$ \\
$A$ & element of order 3 in $\text{Aut}(D4)$; half-integral, hence
  computed exactly via $2A$ \\
$\det(1-A)$ & determinant; $|\det(1-A)| = 9$ counts the orbifold
  fixed points \\
$Z_3$, $C_3$ & cyclic group of order 3 (triality resp.\ fibre
  action) \\
$S_3$, $A_5$ & symmetric group on 3 elements; alternating group on
  5 elements (icosahedral) \\
$\omega$ & third root of unity, $\omega = e^{2\pi i/3}$;
  phase $+120^\circ$ \\
$\chi_0,\chi_1,\chi_2$ & the three $Z_3$ characters; index the
  orbifold sectors \\
\midrule
\multicolumn{2}{l}{\textit{Hilbert space and fibre}} \\
$\mathcal{H}$ & state space,
  $\mathcal{H} = L^2(\text{T}^4)\otimes\mathbb{C}^3$
  (Docs.~230/282) \\
$\mathbb{C}^3$ & fibre over each point; carries the $Z_3$
  structure \\
$3$, $3'$ & the two triplets of the doubling $6 = 3\oplus3'$
  (Doc.~285) \\
$\varphi$ & golden ratio, $\varphi = (1+\sqrt5)/2$; occurs only in
  the $A_5$ layer \\
$\theta = 2/9$ & phase share from Doc.~293; not a lattice invariant
  (Ch.~F) \\
\midrule
\multicolumn{2}{l}{\textit{FFGFT references}} \\
$\xi = 4/30000$ & geometric fundamental constant of FFGFT \\
$\tilde T\cdot m = 1$ & time-mass relation; readable in the
  spectrum as the splitting at $r \ne 1$ \\
\midrule
\multicolumn{2}{l}{\textit{Status labels}} \\
{}[K] & core derivation: computed and controlled \\
{}[S] & structural statement: consistent but not derived \\
theorem & mathematically proven, not merely checked numerically \\
boundary & explicitly booked limitation of scope \\
\end{longtable}
}
\end{center}

% ============================================================
\chapter{The starting position: flatness is not the difference}
% ============================================================

The torus T$^4 = \mathbb{R}^4/\Gamma$ is \emph{always} flat,
whatever lattice $\Gamma$ defines it. The D4 torus and the Z$^4$
torus both have zero curvature. The difference lies not in curvature
but in lattice structure: kissing number $24$ against $8$, packing
density $\pi^2/16$ against $\pi^2/32$.

On passing to $\mathcal{H} = L^2(\text{T}^4)$ the lattice becomes
the \textbf{index set of the Fourier modes}: the basis consists of
plane waves $e^{ik\cdot x}$, and the admissible $k$ lie on the dual
lattice $\Gamma^{*}$. In this sense every lattice becomes flat in
Hilbert space -- it is no longer a sphere packing in space but a
point lattice of quantum numbers.

\textbf{The geometry is not lost; it migrates.} The Laplacian has
eigenvalues $|k|^2$, and the degeneracy of each eigenvalue is the
number of lattice vectors of that length -- the theta series of the
lattice. The 24 shortest D4 vectors become the 24-fold degeneracy of
the first excitation, and the decomposition $24 = 12+12$ of Doc.~311
(FCC neighbours plus winding neighbours) remains exactly readable as
a $k_4$ split.

Two anchors keep this translation stable:

\begin{enumerate}
  \item \textbf{D4 is self-dual up to scaling} -- $D4^{*}$ is again a
        D4 lattice. The Fourier transform does not lead out of the
        family. (Z$^4$ is exactly self-dual.)
  \item \textbf{The fibre changes nothing about the lattice.} In
        $\mathcal{H} = L^2(\text{T}^4)\otimes\mathbb{C}^3$ the
        $\mathbb{C}^3$ fibre hangs at every point and carries the
        $Z_3$ structure; the mode lattice remains $\Gamma^{*}$.
\end{enumerate}

\begin{equation}
  \boxed{\;\text{geometry}\;\longrightarrow\;
  \text{spectrum}\;+\;\text{fibre}\;}
\end{equation}

% ============================================================
\chapter{The spectrum separates what curvature does not [K]}
% ============================================================

Complete shell count up to $|k|^2 = 12$:

\begin{center}
\begin{tabular}{rrrc}
\toprule
Norm $|k|^2$ & Z$^4$ & D4 & D4 split ($k_4{=}0$ / $k_4{\ne}0$) \\
\midrule
$1$ & $8$ & $0$ & --- \\
$2$ & $24$ & $24$ & $12$ / $12$ \\
$3$ & $32$ & $0$ & --- \\
$4$ & $24$ & $24$ & $6$ / $18$ \\
$5$ & $48$ & $0$ & --- \\
$6$ & $96$ & $96$ & $24$ / $72$ \\
$7$ & $64$ & $0$ & --- \\
$8$ & $24$ & $24$ & $12$ / $12$ \\
$9$ & $104$ & $0$ & --- \\
$10$ & $144$ & $144$ & $24$ / $120$ \\
$11$ & $96$ & $0$ & --- \\
$12$ & $96$ & $96$ & $8$ / $88$ \\
\bottomrule
\end{tabular}
\end{center}

Three readings:

\textbf{(i)~Two flat tori are spectrally distinguishable.} Z$^4$
starts at norm~1 with degeneracy~8, D4 at norm~2 with degeneracy~24.
For D4 the odd norms are \emph{entirely} absent -- the parity
condition $\sum k_i$ even excludes them. What curvature does not
distinguish, the spectrum distinguishes at once.

\textbf{(ii)~Independent control passed.} The D4 series
$24, 24, 96, 24, 144, 96 \dots$ is the known theta series
$\theta_{D4} = (\theta_3^4 + \theta_4^4)/2$; the script reproduces
it coefficient by coefficient.

\textbf{(iii)~The $12+12$ decomposition is spectrally
addressable.} The first excitation carries exactly twelve modes with
$k_4 = 0$ (the spatial FCC half) and twelve with $k_4 \ne 0$ (the
winding half). On higher shells the split is asymmetric ($6/18$ at
norm~4, $8/88$ at norm~12): \textbf{the $12/12$ symmetry is a
property of the first shell, not a general theorem.} This must be
booked explicitly -- Doc.~311 argues with $24 = 12+12$, and that
statement holds exactly where it is made.

% ============================================================
\chapter{Anisotropy: where the time-mass relation becomes readable [K]}
% ============================================================

If the fourth direction is compactified with a different radius
($r = R_4/R_{123}$, eigenvalue
$E(k) = k_1^2+k_2^2+k_3^2+(k_4/r)^2$), the first shell splits:

\begin{center}
\begin{tabular}{rll}
\toprule
$r$ & $E$ (winding, 12 modes) & $E$ (FCC, 12 modes) \\
\midrule
$1.0$ & $2.000$ & $2.000$ (degenerate: 24) \\
$1.1$ & $1.826$ & $2.000$ \\
$1.5$ & $1.444$ & $2.000$ \\
$2.0$ & $1.250$ & $2.000$ \\
\bottomrule
\end{tabular}
\end{center}

The FCC modes stay fixed at $E = 2$, the winding modes follow
exactly
\begin{equation}
  E_\text{winding} = 1 + \frac{1}{r^2},
  \qquad
  \Delta = 1 - \frac{1}{r^2}.
\end{equation}
The splitting is a direct measure of the radius ratio -- \textbf{the
place where the $\tilde T\cdot m = 1$ structure becomes readable in
the spectrum}: mass direction and spatial directions carry different
excitation energies as soon as they have different lengths.

\section{The crossing map is discrete and short}

All levels have the form $E = a + b/r^2$ with integer $a, b$. In the
range $1 < r \le 2.5$, among all levels up to norm~8 there are
\textbf{exactly two} crossing radii:

\begin{center}
\begin{tabular}{ll}
\toprule
$r = \sqrt{2}$ & pure winding $(0,0,0,\pm2)$ crosses the FCC-12 \\
$r = \sqrt{3}$ & pure winding crosses the winding-12 \\
\bottomrule
\end{tabular}
\end{center}

The level sequence therefore does not change continuously and
chaotically but at two distinguished root radii. Any ordering
statement about the lowest excitations falls into exactly three
regimes: $r < \sqrt2$, $\sqrt2 < r < \sqrt3$, $r > \sqrt3$. Beyond
$\sqrt3$ the \emph{ground tone} of the spectrum is a pure winding
(at $r = 2$: $E = 1.00$ with 2 modes before $1.25$ with 12 before
$2.00$ with 12).

% ============================================================
\chapter{Time from the mass circle: the reciprocity is mode-wise exact [K]}
% ============================================================

The fourth direction is the time--mass dual one; unrolling is onto
time. Two check questions: how do time periods arise -- and does
$\tilde T\cdot m = 1$ hold only on average, when mass and time
\enquote{distribute equally}?

\section{Time periods as covering-space bookkeeping}

Unrolling the compact fourth direction is the universal covering
$\mathbb{R}\to S^1$; the deck transformation group is
$\mathbb{Z}$, and the unrolled coordinate counts cycle
traversals. A $k_4$ mode $e^{ik_4x_4/R_4}$ thereby becomes a
periodic function with period $T_k = 2\pi R_4/k_4$. Since
$\pi_1(S^1) = \mathbb{Z}$ is abelian, this bookkeeping is
\emph{lossless} -- the same construction that raises the
faithfulness question for non-abelian fundamental groups is
faithful here by structure.

\section{Mode-wise exactness -- and what exactly is exact}

In the deformation spectrum (Ch.~C),
$E(k) = k_1^2+k_2^2+k_3^2 + (k_4/r)^2$ is literally
$E^2 = p_\text{space}^2 + m^2$ with $m_k = k_4/R_4$ --
\textbf{the winding number is the mass number}. Hence
\begin{equation}
  T_k \cdot m_k \;=\; \frac{2\pi R_4}{k_4}\cdot\frac{k_4}{R_4}
  \;=\; 2\pi
  \qquad\text{exactly, for every mode.}
\end{equation}
The reciprocity therefore does \emph{not} hold only on average --
it is satisfied mode-wise identically, because period and mass are
built from the same quantum number $k_4$ and the same radius. The
next section makes precise, however, \emph{which} quantity is exact
here: $T_k$ is first of all an arc length on the circle, i.e.\ a
place interval. Side
result: the $12+12$ splitting (Ch.~C) is, in this reading, the
separation massless ($k_4 = 0$, FCC) versus massive (winding), and
the level crossing at $r > \sqrt3$ means: under strong anisotropy
a pure mass mode becomes the ground tone.

\section{The period is a place first, not a duration}

The quantity $T_k = 2\pi R_4/k_4$ is an \emph{arc length on the
circle}: the distance in the coordinate $x_4$ after which the mode
repeats. It is therefore a \textbf{place interval}, not elapsed
time. Properly it is called $\lambda_4$, and the relation
\begin{equation}
  \lambda_4 \cdot m \;=\; \frac{2\pi R_4}{k_4}\cdot\frac{k_4}{R_4}
  \;=\; 2\pi
\end{equation}
is the \textbf{de Broglie relation} along the mass direction --
wavelength times wavenumber -- not yet a statement about time. It
is exact and always valid, but it is a statement about $x_4$.

\textbf{Where, then, does duration come from?} Only from the
unrolled coordinate. A point of unrolled time is a pair
\begin{equation}
  t \;=\; w\cdot 2\pi R_4 \;+\; x_4,
  \qquad w \in \mathbb{Z} \;\;(\text{winding}),\;\;
  x_4 \in [0, 2\pi R_4),
\end{equation}
that is, \emph{number of revolutions} plus \emph{place on the
cycle}. Duration resides entirely in the counter $w$ -- the deck
transformation of the covering -- and place in the remainder. On
the cycle itself there is no distinguished beginning; every
\enquote{moment} is a place that recurs (cf.\ Doc.~313).

\section{The temporal reading carries a \texorpdfstring{$\gamma$}{gamma} factor}

If one asks for the \emph{oscillation period in unrolled time},
the relevant quantity is $T_\text{osc} = 2\pi/E$ with
$E = \sqrt{p_\text{space}^2 + m^2}$. Then
\begin{equation}
  T_\text{osc}\cdot m \;=\; 2\pi\,\frac{m}{E}
  \;=\; \frac{2\pi}{\gamma}
  \;\le\; 2\pi,
  \qquad\text{equality} \iff p_\text{space} = 0 .
\end{equation}

\begin{center}
\begin{tabular}{rrrr}
\toprule
$p_\text{space}/m$ & $\gamma$ & $T_\text{osc}m/2\pi$ &
  $\lambda_4 m/2\pi$ \\
\midrule
$0$   & $1.00000$ & $1.0000000$ & $1.0000000$ \\
$0.5$ & $1.11803$ & $0.8944272$ & $1.0000000$ \\
$1$   & $1.41421$ & $0.7071068$ & $1.0000000$ \\
$2$   & $2.23607$ & $0.4472136$ & $1.0000000$ \\
$5$   & $5.09902$ & $0.1961161$ & $1.0000000$ \\
\bottomrule
\end{tabular}
\end{center}

\textbf{Place and duration are not the same here.} Exactly $2\pi$
is the \emph{place} reading; the \emph{time} reading carries the
factor $1/\gamma$ and coincides with it only for the mode at rest.
Carrying $T_k$ as a \enquote{time period} and booking
$T_k\cdot m = 2\pi$ as temporal reciprocity would accordingly
confuse the two.

\textbf{Three statements in three directions:}
\begin{center}
\begin{tabular}{lll}
\toprule
Reading & Relation & Equality for \\
\midrule
place ($\lambda_4$, de Broglie) & $=2\pi$ & always \\
time, moving mode & $=2\pi/\gamma \le 2\pi$ &
  $p_\text{space} = 0$ \\
mass mixture (Jensen) & $\ge 2\pi$ & sharp shell \\
\bottomrule
\end{tabular}
\end{center}

\noindent The exact relation $\tilde T\cdot m = $ const therefore
holds for the \textbf{resting, sharp} mode -- there and only there
both correction factors vanish simultaneously. Motion pushes the
product down, mass spread lifts it up. This is no contradiction to
Doc.~306 but the Kaluza--Klein reading of its validity
conditions.

\section{For mixtures the equation becomes an inequality}

If a state superposes several $k_4$ shells with weights $p_k$,
then $\langle T\rangle\langle m\rangle =
2\pi\,\langle k\rangle\langle 1/k\rangle$, and by Jensen
\begin{equation}
  \langle k\rangle\langle 1/k\rangle \;\ge\; 1,
  \qquad\text{equality} \iff \text{exactly one shell occupied.}
\end{equation}
Numerically (script~4): sharp shell (any $k$) $\to$ exactly
$1.0000000000$; $k = 1,2$ equally occupied $\to 1.125$; uniform
$k = 1\ldots50 \to 2.295$; thermal
$p_k \propto e^{-k/T} \to 1.0000227$ ($T = 0.1$) up to
$2.600$ ($T = 10$); 2000 random distributions: minimum $1.077$,
never below~1.

\textbf{Booking -- the conjecture inverts:} the condition for
$\tilde T\cdot m = 1$ is \emph{not} equal distribution but
\textbf{sharpness}: exactly one occupied mass shell. Equal
distribution is, on the contrary, a state with a pronounced
reciprocity excess ($2.29$ for $k \le 50$). The excess
$\langle k\rangle\langle 1/k\rangle - 1$ measures the mass
spread of the state (for small spread
$\approx \mathrm{Var}(k)/\langle k\rangle^2$): a mixture has
neither a sharp period nor a sharp mass, and the product of the
means lies strictly above $2\pi$.

\textbf{Together with the $\gamma$ factor} the full condition
follows: $\tilde T\cdot m$ is constant exactly when the mode
\emph{rests} (no $\gamma$) \emph{and} is sharp (no Jensen excess).
Both deviations have different signs and different causes --
motion in space versus spread in mass.

\textbf{Epistemic status:} lattice-side consistency check of the
mode structure (Kaluza--Klein reading of the fourth direction);
the native FFGFT formulation of the reciprocity is Doc.~306 --
nothing is newly derived here; compatibility is checked: passed,
with the sharpening \enquote{exact per mode, $\ge$ for
mixtures}.

\section{The closure fork is spectrally invisible}

Doc.~313 (after Doc.~295) poses the closure question as a fork:
the defect $d = 100\,\xi = 1/75$ per revolution admits three cases
for the \emph{unrolled} time. With the cycle length
$\tau_c = 2\pi/H_0 = 91.9$~Gyr (Doc.~312):

\begin{center}
\begin{tabular}{llll}
\toprule
Case & Condition & Closure & Period \\
\midrule
A & $d = 0$ & after 1 revolution & $91.9$~Gyr \\
B & $\xi$ frozen & after 75 revolutions & $6892$~Gyr \\
C & $\xi$ running (recursion) & never & log spiral \\
\bottomrule
\end{tabular}
\end{center}

\noindent Case~B follows from the rational rotation number
$1-d = 74/75$: since $\gcd(74,75) = 1$, the trajectory closes
after exactly $75$ revolutions, i.e.\ $75\,\tau_c = 6892.5$~Gyr.
Case~C arises when $\xi$ runs through the recursion; the
accumulated defect becomes logarithmic and the winding turns into
an equiangular spiral.

\textbf{This table is only quoted here.} It stands so that the
cases are intelligible without looking them up; its
\emph{interpretation} -- entropy obligation, Poincaré discrepancy,
thermal history -- belongs to Docs.~312/313 and is not pursued
here. Two points are nevertheless needed so the numbers are not
misread.

\textbf{First: these are deadlines, not repetitions.} $91.9$ and
$6892$~Gyr do not say that something repeats, but within what
period a recurrence would \emph{have to} occur if the respective
case held. What is closed is the \emph{geometry}; whether the
\emph{history} returns is independent of that and, per Doc.~313,
coarse-grained open. In this reading the \enquote{Big Bang} is in
any case a \emph{place} on the cycle -- the region of smallest
masses and slowest clocks -- not an event; a closed circle has no
starting point.

\textbf{Second, and this bounds the reach of all time statements:
states at other places on the cycle are not derivable from
today's.} The usual backward extrapolation ($T(z) = T_0(1+z)$ from
an adiabatic chain) presupposes \emph{one} continuous
thermodynamic history. In the place reading, an early epoch is not
an earlier state of the same system but a location with a
\emph{different local structure} -- different masses, different
clock rates -- and hence its own equilibrium, fixed locally and
not by continuation of today's. The corpus demonstrates this on
itself: $T_0$ is not extrapolated but obtained structurally from
$\xi$ ($k_BT_0 = (16/9)\xi$, Docs.~061/313). This is precisely why
the coarse-grained chaining remains open (Poincaré discrepancy
$10^{104}$~dex, D(ii) in Doc.~313).

\noindent \textbf{What this does \emph{not} mean:} it devalues no
structural derivation. The $\Omega_m^{*}$ chain in Doc.~313 runs
through $\xi$ and $K_\text{frak}$, not through an adiabatic
back-calculation, and is unaffected by this boundary. Affected are
only statements of the form \enquote{how hot was it in epoch~X}
insofar as they are obtained by extending today's equilibrium --
such statements would need their own structural derivation
\emph{at that place}.

It would be tempting to read the defect as a \emph{twist} of the
boundary condition ($k_4 \to k_4 + 1/75$, Scherk--Schwarz-like).
That is \textbf{wrong}, and the reason is structural:

\textbf{In the rolled-up state there are no fractal corrections.}
Single-valuedness of the mode on the circle forces
$k_4 \in \mathbb{Z}$ -- integer or nothing. A winding number is
\emph{topological}; it has no room for $1/75$. Corrections require
something that \emph{accumulates}, and only a traversed path can
accumulate. The very phrase \enquote{defect per revolution}
presupposes counting revolutions -- and revolutions are counted
only in the unrolled coordinate $w$ (Ch.~D1). The fractal path
correction $K_\text{frak}$ is a property of the \emph{traversal},
not of the periodicity of the field.

\textbf{Hence the assignment:}

\begin{center}
\begin{tabular}{lll}
\toprule
Quantity & State & fractal correction \\
\midrule
theta series, degeneracies, $24 = 12+12$ & rolled up & no \\
$\text{Aut}(D4)$, triality, orbifold & rolled up & no \\
circulant, containers $4\times6$ & rolled up & no \\
Casimir, heat kernel, congruences & rolled up & no \\
$\lambda_4\cdot m = 2\pi$ (de Broglie) & rolled up & no \\
\midrule
cycle length, defect $d$, case A/B/C & unrolled & yes \\
revolution count $w$, history & unrolled & yes \\
$T_\text{osc}$, $\gamma$ factor & unrolled & yes \\
\bottomrule
\end{tabular}
\end{center}

\textbf{Two consequences.} First: \emph{the fork has no spectral
signature.} The mode spectrum is an object of the rolled-up state;
whether the trajectory closes after one, after 75 or after no
revolution changes nothing about the eigenvalues of the Laplacian
on T$^4$. The first shell carries 24 modes in all three cases.

Second, and this is the real relief: \textbf{Chapters E--H are
case-independent.} Orbifold, circulant and containers do not
depend on which of the three cases holds -- they are statements
about the compact geometry, and the defect does not act there. The
fork is a question of \emph{history} (entropy, recurrence --
Doc.~313), not of \emph{structure}.

\textbf{The same boundary settles the fractal question.} $D_f$ and
$K_\text{frak}$ are \emph{different} things (Doc.~133):
$D_f^\text{space} = 3-\xi$ is local and practically without effect
($6.7\times10^{-5}$), $K_\text{frak} = 1-100\xi$ is cumulative
across the RG run and acts only where scales are \emph{traversed}.
The rolled-up spectrum carries neither appreciably -- Ch.~J
elaborates.

\textbf{Marginal note on the number.} $75 = 1/(100\xi)$; its
factorisation $3\cdot5^2$ carries \emph{no} meaning -- the five
here has nothing to do with the fivefold symmetry of Ch.~F.

% ============================================================
\chapter{Reorganisations: three levels and one boundary}
% ============================================================

Flat lattices can be reorganised -- on three levels that cost very
different things.

\textbf{Change of basis (free).} Any unimodular transformation from
$GL(4,\mathbb{Z})$ describes the same point lattice with different
generators. The spectrum does not change at all. Pure bookkeeping.

\textbf{Lattice symmetries (free, structure-bearing).}
Reorganisations mapping the lattice onto itself -- this is where the
central finding sits, Ch.~E.

\textbf{Rebuilding into another lattice (costs physics).} Discretely
via glue vectors along the chain
\begin{equation}
  D4 \;\subset\; \text{Z}^4 \;\subset\; D4^{*}
  \qquad\text{(index 2 each)},
\end{equation}
or continuously by deformation in the moduli space
$GL(4,\mathbb{R})/(O(4)\times GL(4,\mathbb{Z}))$. Each step shifts
degeneracies; D4 is a distinguished point there (kissing number 24,
provably unbeatable in 4D).

\textbf{The boundary of the spectral statement.} From dimension~4 on
there exist pairs of flat tori that are \emph{isospectral but not
isometric} (Schiemann). The general theorem \enquote{the spectrum
determines the lattice} is \emph{false} in 4D. Uncritical for D4
itself -- kissing number 24 characterises it uniquely -- but it must
not be booked as a rule.

% ============================================================
\chapter{Triality is built in, not attached [K]}
% ============================================================

\section{Computed on the lattice}

\begin{center}
\begin{tabular}{lr}
\toprule
$|\text{Aut}(D4)|$ & $1152 = |W(F_4)|$ \\
$|\text{Aut}(\text{Z}^4)|$ & $384 = 2^4\cdot4!$ \\
$|\text{Aut}(D4)|/|W(D4)|$ & $1152/192 = 6 = |S_3|$ \\
\bottomrule
\end{tabular}
\end{center}

Both orders were determined by brute-force enumeration of all basis
tuples with identical Gram matrix and match the known values. The
quotient $6 = |S_3|$ establishes triality \emph{directly on the
lattice automorphism group}, not merely on the Dynkin diagram: D4 is
the only lattice of the $D_n$ family with three equivalent Dynkin
arms; for $D5, D6, \dots$ the outer symmetry shrinks to $Z_2$.

\section{The 80 order-3 elements in three classes}

Computed exactly (the triality elements are \emph{half-integral} --
the glue-vector mixture $(\frac12,\frac12,\frac12,\frac12)$ enters;
rounding to integers destroys them, hence the computation uses
$2A$):

\begin{center}
\begin{tabular}{lrl}
\toprule
Class & Count & Character \\
\midrule
trace $-2$, half-integral & $16$ & fixed-point free in
  $\mathbb{R}^4$; eigenvalues $\{\omega,\bar\omega\}$ twice \\
trace $1$, integral & $32$ & fixed plane in $\mathbb{R}^4$,
  free on the 24 minimal vectors \\
trace $1$, half-integral & $32$ & 6 fixed vectors $+$ 6 triple
  orbits \\
\bottomrule
\end{tabular}
\end{center}

\section{Two consequences}

\textbf{(a)~The orbifold is a lattice symmetry, not an extra
assumption.} For each of the 16 elements with trace $-2$ one has
exactly
\begin{equation}
  |\det(1-A)| = 9,
\end{equation}
so the quotient T$^4/\langle A\rangle$ is a $Z_3$ orbifold with
\emph{exactly 9 fixed points} -- the standard number $3^2$. All nine
lie at third-coordinates. \textbf{The T$^4/Z_3$ structure need not
be imposed; it is already present in D4.}

\textbf{(b)~On every mode orbit the $Z_3$ acts as a circulant.} The
unitary action on $L^2(\text{T}^4)$ permutes the three plane waves
of a triple orbit cyclically; the permutation matrix has eigenvalues
$\{1,\omega,\omega^2\}$, phases $0^\circ$, $+120^\circ$,
$-120^\circ$ -- \emph{structurally exactly the circulant
diagonalisation of the $\mathbb{C}^3$ fibre} (Doc.~291). The first
shell provides eight copies: $24 = 8\times3
= 8\times(\chi_0\oplus\chi_1\oplus\chi_2)$.

% ============================================================
\chapter{The negative finding: 5 does not divide 1152 [K]}
% ============================================================

The order distribution in $\text{Aut}(D4)$, determined exactly:

\begin{center}
\begin{tabular}{lrrrrrrr}
\toprule
Order & $1$ & $2$ & $3$ & $4$ & $6$ & $8$ & $12$ \\
\midrule
Count & $1$ & $139$ & $80$ & $228$ & $464$ & $144$ & $96$ \\
\bottomrule
\end{tabular}
\end{center}

Total $1152$. \textbf{Order 5: zero elements} -- and that is not a
computational result but a theorem: $1152 = 2^7\cdot3^2$, and $5$
does not divide it (Lagrange).

\textbf{Consequence for $\theta = 2/9$.} The fivefold rotation
$R_5$, from which $p_0 = |\langle v_0|R_5 v_e\rangle|^2 = 2/9$
follows (Doc.~293), is \emph{in principle} incompatible with the D4
lattice structure. The test was additionally carried out directly:
all invariants of the linear lattice action at the nine orbifold
fixed points ($x\!\cdot\!x$, $k\!\cdot\!x$, $x\!\cdot\!y$, all mod 1
in exact rational arithmetic) have denominator spectrum
$\{1,2,3,6\}$ -- \textbf{no ninths}. Notably, $x\!\cdot\!y$ is a
multiple of one ninth, but the numerators are always divisible by 3;
the D4 geometry \emph{suppresses} ninths at the first level.

\textbf{What this means -- and what it does not.} $2/9$ is not an
invariant of the lattice action. This does not devalue Doc.~293 but
delimits it: the phase belongs to a layer \emph{above} the lattice
(the $A_5$ structure), the lattice carries the $Z_3$. The division
of labour is thereby exactly named rather than surmised.

% ============================================================
\chapter{The doubling is nevertheless prepared [K]}
% ============================================================

Does the negative finding contradict the doubling $6 = 3\oplus3'$ of
Doc.~285? No -- in three specifiable stages.

\textbf{(1)~The prohibition hits the wrong level.} $5\nmid1152$
forbids order-5 \emph{lattice automorphisms}. The doubling lives on
the \textbf{fibre} ($\mathbb{C}^3\to\mathbb{C}^6$); $A_5$ is nowhere
required as a spatial symmetry. The compatibility condition is only
that the interface of the two levels exists -- and that is exactly
determinable: element orders in $A_5$ are $\{1,2,3,5\}$, in
$\text{Aut}(D4)$ they are $\{1,2,3,4,6,8,12\}$; the intersection is
$\{1,2,3\}$, canonically $C_3$. \textbf{And exactly that is provided
by triality.} The interface is not merely permitted, it is the only
possible one -- and it is occupied.

\textbf{(2)~The lattice prepares the doubling.} The central
involution $-1 \in \text{Aut}(D4)$ commutes with the triality $Z_3$
and pairs the eight triple orbits of the first shell into
\textbf{four disjoint pairs}. No orbit is self-paired, provably:
$-k = A^j k$ would require eigenvalue $-1$ for an element of
order~3 -- impossible. Hence
\begin{equation}
  24 \;=\; 8\times3 \;=\; 4\times6
  \qquad\text{(orbit $\oplus$ mirror orbit).}
\end{equation}
Each six-block carries under $C_3$ the character content
$2\times(\chi_0\oplus\chi_1\oplus\chi_2)$ -- exactly the restriction
of $3\oplus3'$. The first excitation shell provides four ready
containers.

\textbf{(3)~The boundary, cleanly booked.} The \emph{labelling} --
which factor is $3$ and which $3'$ -- is lattice-blind. The $A_5$
characters of the two triplets differ only on the 5-classes
($\varphi$ against $1-\varphi$); on the 3-class both are zero and
the $C_3$ restrictions are identical. The distinction hangs exactly
on the structure that does not exist in the lattice -- fully
consistent with Ch.~F.

\textbf{Balance:} the lattice provides containers and $C_3$ content,
the $A_5$ layer provides the $3/3'$ labels, $\varphi$ and the
fivefold -- and hence $2/9$.

% ============================================================
\chapter{Anchor-free results [K]}
% ============================================================

All quantities below are dimensionless -- ratios, counts,
congruences. No Kammerton, no SI scale enters.

\section{Two congruence theorems for the D4 theta series}

All complete shells up to norm~120 checked:
\begin{equation}
  N(n) \equiv 0 \pmod 3
  \qquad\text{and}\qquad
  N(n) \equiv 0 \pmod 6
  \qquad\text{without exception.}
\end{equation}
Both are theorems, not empirics: the triality $Z_3$ is fixed-point
free on $\Gamma\setminus\{0\}$ (eigenvalues $\omega,\bar\omega$ --
no eigenvalue 1), so \emph{every} shell decomposes into triple
orbits; together with the $-1$ pairing without self-pairing, into
six-blocks.

\section{Character-resolved T$^4/Z_3$ spectrum}

Consequence: exact equidistribution on every shell.

\begin{center}
\begin{tabular}{rrrrr}
\toprule
Norm & $N(D4)$ & $\chi_0$ & $\chi_1$ & $\chi_2$ \\
\midrule
$2$ & $24$ & $8$ & $8$ & $8$ \\
$4$ & $24$ & $8$ & $8$ & $8$ \\
$6$ & $96$ & $32$ & $32$ & $32$ \\
$10$ & $144$ & $48$ & $48$ & $48$ \\
$14$ & $192$ & $64$ & $64$ & $64$ \\
$18$ & $312$ & $104$ & $104$ & $104$ \\
\bottomrule
\end{tabular}
\end{center}

The invariant sector of the orbifold ($\chi_0$) has exactly $N(n)/3$
states per shell; the $Z_3$ projection costs no asymmetry on any
shell.

\textbf{Honest consequence.} Equidistribution and six-containers are
\emph{not} a distinction of the first shell but a theorem for all.
\textbf{The structure is everywhere -- and therefore cannot select
anything.} This is exactly the pattern-versus-occupation logic of
Doc.~293: the pattern is present scale-invariantly; which shell is
physically occupied is decided by something else.

\section{Casimir: the densest lattice loses}

$\zeta$-regularised vacuum energy of a scalar field,
$E = \pi\,\zeta_M(-1/2)$, both tori normalised to covolume~1
(Epstein zeta via the Riemann split representation, verified three
times -- against direct summation at $s = 4, 5$ and against the
closed form
$Z_{\text{Z}^4}(s) = 8(1-4^{1-s})\zeta(s)\zeta(s-1)$, agreement to
25 digits):

\begin{center}
\begin{tabular}{lrr}
\toprule
 & $\zeta_M(-1/2)$ & $E = \pi\zeta_M(-1/2)$ \\
\midrule
Z$^4$ torus & $-0.2966893$ & $-0.9320768$ \\
D4 torus & $-0.2764778$ & $-0.8685805$ \\
\midrule
Difference & & $+0.0634963$ \\
\bottomrule
\end{tabular}
\end{center}

\textbf{The Z$^4$ torus lies lower.} The reason is the spectral gap:
at equal volume the densest lattice D4 has the \emph{largest} first
eigenfrequency ($\sqrt2$ against $1$) and hence the \emph{less}
negative zero-point energy. Notable is the reversal of ordering: for
$s > 0$ D4 minimises the Epstein zeta (the known extremality of the
densest lattice), at $s = 0$ all lattices agree
($\zeta(0) = -1$ universally), and at $s = -1/2$ the ordering is
inverted.

\textbf{Interpretation boundary, booked:} this holds for the scalar,
periodic field. Different field content or different boundary
conditions can change the sign of the preference --
\enquote{nature chooses Z$^4$} does \emph{not} follow.

\section{Heat kernel: the observability limit quantified}

For both tori $t^2\theta(t) \to 1$ (identical Weyl term at equal
volume); the lattice difference is exponentially small:

\begin{center}
\begin{tabular}{rr}
\toprule
$t$ & relative difference \\
\midrule
$0.5$ & $1.2\times10^{-2}$ \\
$0.2$ & $1.2\times10^{-6}$ \\
$0.1$ & $1.8\times10^{-13}$ \\
$0.05$ & $4.1\times10^{-27}$ \\
\bottomrule
\end{tabular}
\end{center}

\textbf{Any observable resting only on the leading heat-kernel
coefficients is lattice-blind.} The D4/Z$^4$ difference sits
exclusively in the exponentially suppressed theta-series fine
structure -- or directly in the spectral gap.

% ============================================================
\chapter{What is still bare here -- and what is not [K]}
% ============================================================

This document computes with the \emph{smooth} Laplacian on
$\text{T}^4 = \mathbb{R}^4/\Gamma$. What has to be settled is what
FFGFT holds against this in fractal structure -- and here
\textbf{two quantities must be kept apart} that are easily
confused:

\begin{center}
\begin{tabular}{lll}
\toprule
Quantity & Character & Effect \\
\midrule
$D_f^\text{space} = 3-\xi$ (Doc.~133) & \emph{local} &
  $6.67\times10^{-5}$ \\
$K_\text{frak} = 1-100\xi$ (Doc.~133, A040) & \emph{cumulative} &
  $1.33\times10^{-2}$ \\
\bottomrule
\end{tabular}
\end{center}

\noindent Doc.~133 explicitly derives $K_\text{frak}$ \emph{not}
from the local dimension but from the \enquote{cumulative effect}
across the RG run from the Planck to the electroweak scale (17
orders of magnitude); the factor 100 is RG amplification
(Doc.~190). The local dimension alone yields, per Doc.~133,
\enquote{a vanishingly small correction -- practically one},
computed $(D/3)^{D/2} = 0.99993$. The ratio of the two effects is
exactly $200$.

\textbf{On notation, because this is easily slipped:} the fractal
dimension of FFGFT is $D_f^\text{space} = 3-\xi$ -- the
\emph{spatial} one, not a four-dimensional quantity. The exponent
$3/2$ in $K_\text{frak}^{3/2}$ (Doc.~018) is likewise justified as
half the \emph{spatial} dimension. It is three dimensions, not
four; a quantity $D_f = 4-\xi$ is unknown to the corpus.

\textbf{Hence the assignment} -- and it coincides with the
boundary of Ch.~D2 (rolled up versus unrolled): a Laplacian on the
compact T$^4$ is \emph{local}. It cannot carry a cumulative
quantity; $K_\text{frak}$ does not enter through the operator but
through the \textbf{bridge constants} ($E_0$, $v$) in scale
conversion -- exactly as R72 shows for the two $E_0$ values
(ratio $1/\sqrt{K_\text{frak}}$). The consequence falls into
three very unequal parts.

\section{Untouched: everything anchor-free}

By the assignment rule (register~R70, Doc.~313), $K_\text{frak}$
applies to absolute quantities and cancels for ratios of the same
level. Untouched therefore:

\begin{itemize}
  \item all \textbf{counts} -- $24$, $8$, $N(n)$, $8\times3$,
        $4\times6$: integer combinatorics, no length;
  \item the \textbf{congruences} $N(n)\equiv0 \bmod 3$ and
        $\bmod\,6$: group theory, no metric content;
  \item the \textbf{character split} $N/3$ per sector;
  \item all \textbf{ratios} $E/E_\text{min}$ and the crossing radii
        $\sqrt2$, $\sqrt3$ -- structure over structure, the factor
        cancels.
\end{itemize}

\noindent This is the real reason why Ch.~H is called
\enquote{anchor-free results}: those statements are
fractal-invariant.

\section{One factor: all absolute energies}

With an anchor, $E = \hbar H_0\,|k|$ (ground mode $= \hbar H_0$, see
Ch.~B together with Doc.~312). These values are bare and carry the
factor:

\begin{center}
\begin{tabular}{lrr}
\toprule
 & bare & with $K_\text{frak}$ \\
\midrule
first D4 shell ($|k| = \sqrt2$) & $3.230\times10^{-52}$~J &
  $3.187\times10^{-52}$~J \\
pure winding ($|k| = 2$) & $4.567\times10^{-52}$~J &
  $4.507\times10^{-52}$~J \\
\bottomrule
\end{tabular}
\end{center}

\noindent Throughout $-1.33\,\%$. The same holds for the Casimir
energies of Ch.~H and for the spectral gap itself. Bookkeeping,
settled; not structurally interesting.

\textbf{Important for the assignment:} the factor sits \emph{not}
in the spectrum but in the \emph{anchor}. The dimensionless
eigenvalues $|k|^2$ remain unchanged; what is corrected is the
conversion to SI via $\hbar H_0$ -- a scale bridge, that is,
precisely the route along which a cumulative quantity may act.
Anyone looking only at ratios sees none of it (A040 rule, R72).

\section{The operator, not the conversion: the perturbation calculation [K]}

The third part cannot be settled by multiplication. At
$D_f^\text{space} = 3-\xi$ the operator is \emph{no longer an
exactly smooth Laplacian} -- and exact symmetry is precisely what
carries the degeneracies. Only the \emph{local} effect is relevant
here: $K_\text{frak}$ is cumulative per Doc.~133 and therefore
does not belong on a local operator, so $O(100\xi)$ is ruled out
as a perturbation strength.

\begin{center}
\begin{tabular}{lrr}
\toprule
Perturbation strength & $\Delta E/E$ & $\Delta E$ at
  $E = \hbar H_0$ \\
\midrule
local, $(D/3)^{D/2}$ (Doc.~133) & $6.7\times10^{-5}$ &
  $1.5\times10^{-56}$~J \\
$O(\xi)$ (upper bound) & $1.3\times10^{-4}$ &
  $3.0\times10^{-56}$~J \\
\bottomrule
\end{tabular}
\end{center}

\noindent The perturbation thus lies \textbf{two orders of
magnitude below} what an estimate via $K_\text{frak}$ would
give.

\noindent \textbf{The 24-fold degeneracy would then no longer be a
degeneracy but a multiplet with fine structure.} Which
sub-multiplets it decomposes into is decided by the \emph{symmetry
of the perturbation} -- and, as the calculation shows, via
\emph{representation theory}, not via the orbit structure:

\begin{center}
\begin{tabular}{ll}
\toprule
Perturbation respects & Multiplet structure of the 24 (computed) \\
\midrule
full $\text{Aut}(D4)$ ($c_g$ norm-dependent only) &
  $9 + 8 + 4 + 2 + 1$ \\
triality ($Z_3$) and $-1$ & $8\times1 \;+\; 8\times2$ \\
only $Z_3$ & $8\times1 \;+\; 8\times2$ \\
only $-1$ & $24\times1$ \\
neither & $24\times1$ \\
\bottomrule
\end{tabular}
\end{center}

\textbf{Computed} (script~4): first-order perturbation theory on
the $24\times24$ matrix $\langle k|V|k'\rangle = c_{k-k'}$,
symmetry classes via orbit averaging of the coefficients, three
independent random perturbations per class -- the multiplet
structure is seed-independent and identical in all cases.

\textbf{The obvious expectation is wrong.} One assumes the orbits
stay together ($Z_3$ and $-1$: the 24 as a whole; only $Z_3$:
8~triple orbits; only $-1$: 4~six-blocks). In fact multiplet sizes
follow the \emph{irreducible representations} of the respected
symmetry group, not the orbits. Concretely: (i)~every position-dependent
perturbation breaks the translation invariance of the torus --
and \emph{that} carried the 24-fold degeneracy, not the point
symmetry. (ii)~$Z_3$ is abelian, all irreps one-dimensional; the
doublets are the antiunitary pairing of the conjugate characters
$\chi_1\leftrightarrow\chi_2$, while the $\chi_0$ sector
decomposes into 8~singlets. (iii)~$-1$ alone enforces nothing
(free $Z_2$ action, one-dimensional characters). (iv)~Even the
maximally symmetric perturbation does not keep the 24 together:
$24 = 9\oplus8\oplus4\oplus2\oplus1$, consistent with the
irrep dimensions of $W(F_4)$ on the root shell [K numerically;
the irrep assignment is B, without a character computation].

\textbf{This closes the open point -- and defuses it:}
qualitatively it still holds that only as much of the 24-fold
degeneracy survives as the perturbation symmetry provides in irrep
dimension (at best $9+8+4+2+1$, generically nothing).
Quantitatively, however, the relevant perturbation is
$\sim10^{-5}$ relative, not $10^{-2}$: the splitting would be fine
structure far below any readability, and the container structure
($4\times6$, Ch.~G) remains practically intact.

\textbf{Chapter balance.} \emph{Bare} are only the SI absolute
values, and via the anchor, not via the operator. The rolled-up
quantities of this document -- degeneracies, ratios, congruences,
containers, Casimir and heat-kernel ratios -- are \textbf{not}
bare but correct: $K_\text{frak}$ does not act there
structurally. This is the same boundary as in Ch.~D2, read from
the other side.

% ============================================================
\chapter{Status overview}
% ============================================================

{\small
\begin{longtable}{lp{6.4cm}l}
\toprule
Statement & Content & Status \\
\midrule
\endfirsthead
\toprule
Statement & Content & Status \\
\midrule
\endhead
\bottomrule
\endfoot
Translation & geometry $\to$ spectrum $+$ fibre; lattice becomes
  mode index, degeneracies $=$ theta series & [K] \\
Distinguishability & Z$^4$ from norm 1 (8), D4 from norm 2 (24);
  odd norms empty for D4 & [K] \\
$12+12$ & first shell exactly FCC $+$ winding; higher shells
  asymmetric ($6/18$, $8/88$) & [K] \\
Anisotropy & $E_\text{wind} = 1 + 1/r^2$; crossings only at
  $r = \sqrt2$ and $\sqrt3$; three regimes & [K] \\
Triality & $|\text{Aut}(D4)| = 1152$, quotient $6 = |S_3|$;
  80 order-3 elements in three classes & [K] \\
Orbifold & 16 elements with $|\det(1-A)| = 9$: T$^4/Z_3$ is a
  lattice symmetry, not an extra assumption & [K] \\
Circulant & $Z_3$ acts on every triple orbit with
  $\{1,\omega,\omega^2\}$ $=$ $\mathbb{C}^3$ fibre structure & [K] \\
$5 \nmid 1152$ & no order-5 automorphisms; $2/9$ is not a lattice
  invariant (denominators $\{1,2,3,6\}$) & [K] theorem \\
Doubling & interface $A_5 \cap \text{Aut}(D4) = C_3$, occupied;
  $24 = 4\times6$ as containers & [K] \\
Labelling & $3$ versus $3'$ is lattice-blind (difference only on
  5-classes) & [K] boundary \\
Congruences & $N(n) \equiv 0 \bmod 3$ and $\bmod 6$ on all 60
  shells checked, with proof ground & [K] theorem \\
Selection & containers and equidistribution are generic --
  the structure cannot distinguish any shell & [K] \\
Casimir & $E(\text{Z}^4) = -0.9321$ against
  $E(D4) = -0.8686$; ordering reversed relative to $s>0$
  & [K], field-content dependent \\
Heat kernel & Weyl term identical, difference $4\times10^{-27}$ at
  $t = 0.05$: leading order is lattice-blind & [K] \\
Spectral theorem & \enquote{spectrum determines lattice} false in 4D
  (Schiemann); uncritical for D4, but not a rule & [K] boundary \\
Reciprocity & $\lambda_4\cdot m = 2\pi$ exact as a place relation
  (de Broglie); temporally $2\pi/\gamma \le 2\pi$; for mixtures
  $\ge 2\pi$ (Jensen) -- exact only at rest and sharp & [K] \\
Closure case & rolled up, no fractal correction ($k_4$
  topological); fork A/B/C spectrally invisible, Chs.~E--H
  case-independent & [K] \\
Perturbation multiplets & splitting follows irrep dimensions, not
  orbits: $9{+}8{+}4{+}2{+}1$ (full), $8{\times}1{+}8{\times}2$
  ($Z_3$), $24{\times}1$ (generic) & [K] \\
Fractal status & $D_f^\text{space} = 3-\xi$ local
  ($6.7{\times}10^{-5}$) versus $K_\text{frak}$ cumulative
  ($1.33\,\%$, RG run, Doc.~133): bare are only SI absolute values
  via the anchor, not the rolled-up spectrum & [K] \\
\end{longtable}
}

\medskip
\noindent\textbf{Balance:} in Hilbert space the lattice becomes flat
but not inconsequential. The kissing number becomes the degeneracy,
triality becomes the orbifold and the fibre, the radius ratio
becomes the level splitting. The \emph{division of labour} is
determined just as sharply: fivefold symmetry -- and with it
$\varphi$ and $2/9$ -- cannot be carried by the lattice
\emph{itself}, and that by theorem ($5 \nmid 1152$), not for want
of computational effort. This is a statement about the
\emph{layer}, not about the quantity: on the fibre, $2/9$ is
explicitly admissible, since the interface
$A_5 \cap \text{Aut}(D4) = C_3$ is occupied by triality (Ch.~G),
and with $24 = 4\times6$ the lattice even provides the containers
for the doubling $3\oplus3'$. What is excluded is therefore not
$2/9$ but only its derivation \emph{from} lattice invariants --
the phase demonstrably enters one layer higher. And the container
structure is generic: it provides, it does not select.

% ============================================================
\chapter{Check scripts}
% ============================================================

Three scripts, all deterministic, all controls passed:

\begin{itemize}
  \item \texttt{d4\_skript\_1\_spektrum\_deformation.py} --
        theta series D4/Z$^4$, $k_4$ split, deformation $r$, exact
        crossing radii (standard library only).
  \item \texttt{d4\_skript\_2\_trialitaet\_orbifold\_phasen.py} --
        $\text{Aut}(D4) = 1152$, order-3 classification (computed
        exactly with $2A$), $|\det(1-A)| = 9$, circulant, phase test
        $2/9$, order distribution, $24 = 4\times6$ (numpy).
  \item \texttt{d4\_skript\_4\_stoerung\_reziprozitaet.py} --
        first-order perturbation theory on the 24 shell (multiplet
        structure per symmetry class, 3 seeds each) and
        reciprocity $T\cdot m$ (mode-wise; Jensen for mixtures)
        (numpy, mpmath).
  \item \texttt{d4\_skript\_3\_schalen\_casimir.py} --
        shell theorems, Epstein zeta at $s = -1/2$, heat kernel
        (numpy, mpmath).
\end{itemize}

Controls against known values: $\theta_{D4}$, $|W(F_4)| = 1152$,
$2^4\cdot4! = 384$, $|W(D4)| = 192$, orbifold fixed-point count
$3^2 = 9$, Jacobi identity for $Z_{\text{Z}^4}$.

\medskip
\noindent\textbf{Registration:} entry into Doc.~190 will be done
separately, should the delimitation against Doc.~293 be booked
there.

\end{document}
